\documentclass[twocolumn]{aastex631}
\begin{document}
\title{An age dependence of the Galactic alpha-knee across galactocentric radius}
\author{NebulaMind Lab (autonomous pipeline)}
\affiliation{NebulaMind Lab, \url{https://lab.nebulamind.net}}
\begin{abstract}
Age-resolved alpha-knee from an APOGEE (DR18) 3-table join of 330,457 giants (apogeeDistMass C/N ages + distances, apogeeStar Galactic coordinates, aspcapStar [Fe/H]-[MG/Fe]). The [Fe/H]\_knee is located per (R\_g x age) cell with a broken-line ridge fit (bootstrap error) in 41 populated cells; median [Fe/H]\_knee = -0.50. Gradients: d[Fe/H]\_knee/d(age)|\_R\_g = +0.0266+/-0.0105 dex/Gyr; d[Fe/H]\_knee/dR\_g|\_age = -0.0557+/-0.0033 dex/kpc. Non-circularity: the age-gradient sign HOLDS under the abundance-scale swap ([Fe/H]->[M/H], [MG/Fe]->[alpha/M]). Generated autonomously from public data (apogee) via the NebulaMind Lab runner: a bounded, reproducible, descriptive study.
\end{abstract}
\section{Introduction}
Introduction:
The study of the Milky Way's chemical evolution is crucial for understanding its formation and development. Previous works have investigated various aspects of this topic, such as the role of stellar ages in chemical enrichment [Claytor2020] and the identification of young alpha-rich stars in different Galactic regions [Grisoni2024]. However, a comprehensive analysis of the age-resolved alpha-knee is still lacking. This research aims to fill this gap by examining the relationship between the alpha-knee, stellar ages, and Galactocentric radius using data from the APOGEE catalog.
\section{Data and method}
Data and method:
To achieve this goal, we utilized data from three tables in the SDSS DR18 APOGEE catalog: apogeeDistMass, apogeeStar, and aspcapStar. These tables were joined on APOGEE\_ID and filtered using quality cuts to ensure reliable results. We computed R\_g values based on Galactic coordinates and distances, assuming R0=8.122 kpc. The alpha-knee was determined by fitting a broken-line ridge to the [Fe/H] distribution in each (R\_g x age) cell, with bootstrap errors estimated for uncertainty.
\section{Result}
Result:
Our analysis revealed an age-resolved alpha-knee from 330,457 giants in the APOGEE catalog. We found that the median [Fe/H]\_knee value is -0.50. The gradients of the alpha-knee with respect to age and R\_g were determined as d[Fe/H]\_knee/d(age)|\_R\_g = +0.0266+/-0.0105 dex/Gyr and d[Fe/H]\_knee/dR\_g|\_age = -0.0557+/-0.0033 dex/kpc, respectively. Notably, the age-gradient sign remained consistent even when swapping the abundance calibration scale.
\begin{figure}\centering\includegraphics[width=\columnwidth]{result.png}\caption{An age dependence of the Galactic alpha-knee across galactocentric radius}\end{figure}
\section{Caveats}
Caveats:
It is essential to acknowledge that this study relies on an automated, single-selection, uncalibrated measurement of stellar ages and abundances. The use of C/N-based spectroscopic ages may introduce systematic errors due to their data-driven nature. Additionally, our analysis did not account for potential biases in the APOGEE catalog or uncertainties in the distance calculations. These limitations highlight the need for further research and refinement of methods to improve the accuracy and reliability of age-resolved alpha-knee measurements.
\section*{References}
\footnotesize
[Claytor2020] Claytor et al. (2020). Chemical Evolution in the Milky Way: Rotation-based Ages for APOGEE-Kepler Cool Dwarf Stars. bibcode 2020ApJ...888...43C \par [Grisoni2024] Grisoni et al. (2024). K2 results for "young" α-rich stars in the Galaxy. bibcode 2024A\&A...683A.111G \par [Valle2024] Valle et al. (2024). Stellar model tests and age determination for RGB stars from the APO-K2 catalogue. bibcode 2024A\&A...690A.323V \par [Warfield2021] Warfield et al. (2021). An Intermediate-age Alpha-rich Galactic Population in K2. bibcode 2021AJ....161..100W

\end{document}
